\newpage
\phantomsection
\label{prf:fv-msml-diameter-monotone-under-inclusion}
\LRAProofFor{thm:fv-msml-diameter-monotone-under-inclusion}

\begin{remark*}[Return]
Return to \hyperref[thm:fv-msml-diameter-monotone-under-inclusion]{Diameter Is Monotone}.
\end{remark*}

\begin{theorem*}[Diameter Is Monotone]
Let $(X,d)$ be a metric space. If \(A\subseteq B\), the pairwise distance set
for \(A\) is nonempty, and the pairwise distance set for \(B\) is bounded
above, then
\[
  \operatorname{diam}(A)\leq \operatorname{diam}(B).
\]
\end{theorem*}

\begin{proof}[Professional Standard Proof]
\LRAProofBodyStart
Every pairwise distance realized by two points of \(A\) is also realized by
two points of \(B\), because \(A\subseteq B\). Thus the pairwise-distance set
for \(A\) is a subset of the pairwise-distance set for \(B\). Since the former
is nonempty and the latter is bounded above, monotonicity of the supremum gives
\[
  \sup \{d(x,y)\mid x,y\in A\}
  \leq
  \sup \{d(x,y)\mid x,y\in B\}.
\]
By the definition of diameter, this is exactly
\(\operatorname{diam}(A)\leq\operatorname{diam}(B)\).
\end{proof}

\begin{proof}[Detailed Learning Proof]
\LRAProofBodyStart
The Lean proof does not unfold the supremum by hand. It identifies the exact
order-theoretic theorem needed and then supplies its hypotheses.

First, the target
\[
  \operatorname{diam}(A)\leq \operatorname{diam}(B)
\]
means
\[
  \sup D_A
  \leq
  \sup D_B,
\]
where \(D_A\) is the Lean set \texttt{diameterSet A} and \(D_B\) is the Lean
set \texttt{diameterSet B}.
The theorem \texttt{csSup\_le\_csSup} proves this kind of statement. It asks
for three pieces of data:
\[
  D_B\ \text{is bounded above},\qquad
  D_A\ \text{is nonempty},
\]
and
\[
  D_A\subseteq D_B.
\]

The first two are exactly the hypotheses
\texttt{B\_diameterSet\_bddAbove} and
\texttt{A\_diameterSet\_nonempty}. The only constructed ingredient is the
third one. To construct it, take an arbitrary real number \(r\) in
\texttt{diameterSet A}. By the definition of \texttt{diameterSet}, this means
there are witnesses \(x,y\in A\) with
\[
  r=d(x,y).
\]
Using the inclusion \(A\subseteq B\), the same witnesses \(x\) and \(y\) also
lie in \(B\). Therefore the same equality \(r=d(x,y)\) proves
\[
  r\in D_B.
\]
This witness conversion is packaged in the Lean helper
\texttt{diameterSet\_mono set\_inclusion}.

After those three obligations are supplied, \texttt{csSup\_le\_csSup}
returns the desired inequality directly.
\end{proof}

\begin{proof}[Formalized Lean Proof]
\begin{verbatim}
theorem diameter_monotone_under_inclusion
    {X : Type u}
    [MetricSpace X]
    {A B : Set X}
    (set_inclusion : A subset B)
    (A_diameterSet_nonempty : (diameterSet A).Nonempty)
    (B_diameterSet_bddAbove : BddAbove (diameterSet B)) :
    diameter A <= diameter B := by
  exact csSup_le_csSup
    B_diameterSet_bddAbove
    A_diameterSet_nonempty
    (diameterSet_mono set_inclusion)
\end{verbatim}
\end{proof}

\begin{remark*}[Proof structure]
The helper lemma \texttt{diameterSet\_mono} performs the metric-space-specific
work. The final theorem is then an order-theoretic application of
\texttt{csSup\_le\_csSup}.
\end{remark*}

\NoLocalDependencies

\clearpage
